% ============================================================================
\chapter{Ступень VI. Биология, генетический код ДНК и морфогенез органов}
\label{ch:course-stage6-biology}
% ============================================================================

\section{Урок 19. Делитель водородной связи \texorpdfstring{$N_{\mathrm{HB}} = 63$}{N\_HB = 63} и три планетарных окна жизни (\texorpdfstring{$7 : 2 : 1$}{7 : 2 : 1})}

Перейдём от неживых кристаллов к живой клетке. Жизнь работает в жидкой воде при комнатной температуре, где тепловой фон многократно меньше энергии ковалентной связи, но соизмерим с энергией \textbf{водородной связи}. Чему равен шаг водородной связи в нашем носителе? Потолок амплитуды $K_P(3) = 64$ без вакуумного пола $\mu_{\mathrm{vac}} = 1$ даёт:
\[
  N_{\mathrm{HB}}(3) \;=\; K_P(3) - \mu_{\mathrm{vac}}(3) \;=\; 64 - 1 \;=\; \mathbf{63 \;=\; 7 \times 9 \;=\; (2^3 - 1)\cdot d^2}.
\]
Отсюда сразу следуют три фундаментальных биологических инварианта:
\begin{enumerate}[label=\textbf{\arabic*.},noitemsep]
\item \textbf{Три планетарных фотонных окна биосферы ($441 : 126 : 63 = 7 : 2 : 1$) и оптическое окно фотосинтеза $\bigl[\frac{126}{13},\ \frac{63}{4}\bigr]$:}
  для трёх допустимых режимов среды $m_{\mathrm{env}}\in\{-1, 0, +1\}$ безразмерный приведённый шаг мягкого фотона $\widehat{\Lambda}_{\mathrm{soft}}^{(m)} = \bar{\lambda}_{\mathrm{soft}}^{(m)}/(a_0\alpha_{\mathrm{fs}}^{-1})$ кратен $N_{\mathrm{HB}} = 63$ с коэффициентами $\{d_{\mathrm{out}}(V),\,(c_0/c)^2,\,\mu_{\mathrm{vac}}\} = \{\mathbf{7, 2, 1}\}$, а по единственному ридберговскому проектору $\lambda_{\mathrm{SI}}(\widehat{\Lambda}) = \widehat{\Lambda}\cdot(2\pi a_0\alpha_{\mathrm{fs}}^{-1})$ ($2\pi a_0\alpha_{\mathrm{fs}}^{-1} \approx 45{,}56335\text{ нм}$) это даёт:
  \begin{itemize}[noitemsep]
  \item инфракрасные тепловые окна растворителя трёх миров $\widehat{\Lambda}_{\mathrm{soft}}^{(-1,0,+1)} = (441,\ 126,\ 63) \longmapsto \lambda_{\mathrm{SI}}^{(-1,0,+1)} \approx \mathbf{(20{,}09,\ 5{,}741,\ 2{,}870)\text{ мкм}}$;
  \item оптическое окно земного фотосинтеза ($m_{\mathrm{env}} = 0$, когерентное поглощение от $|{\triangle}(3)| = 8$ до $N_\varphi(3) = 13$ квантов водородной связи $E_{\mathrm{HB}}$): безразмерные шаги $\widehat{\Lambda}_{13}^{(0)} = \frac{126}{13}$ и $\widehat{\Lambda}_8^{(0)} = \frac{126}{8} = \frac{63}{4}$ через тот же ридберговский множитель $45{,}56335\text{ нм}$ дают синий пик Соре хлорофилла $\frac{126}{13}\times 45{,}56335\text{ нм} \approx \mathbf{441{,}61\text{ нм}}$ и красный пик фотосистемы $\frac{63}{4}\times 45{,}56335\text{ нм} \approx \mathbf{717{,}62\text{ нм}}$!
  \end{itemize}
\item \textbf{Точные рациональные критические показатели 3D перехода жидкость--пар (класс Изинга):}
  на оболочках кубооктаэдра критические индексы выражаются точными дробями:
  \[
    \nu = \frac{17}{27},\quad \alpha = \frac{1}{9},\quad \eta = \frac{2}{55},\quad \beta = \frac{323}{990},\quad \gamma = \frac{68}{55},\quad \delta = \frac{91}{19},
  \]
  которые \textbf{тождественно} удовлетворяют всем скейлинговым законам Рашбрука--Уидома--Фишера--Гиперскейлинга ($\alpha + 2\beta + \gamma = \frac{110 + 646 + 1224}{990} = \frac{1980}{990} = 2$ ровно; $2 - \alpha = \frac{17}{9} = 3\nu$ ровно!).
\item \textbf{Почему жизнь зарождается САМА: сквозной изоморфизм второй сферы $S_2 = 42$ и автоматический отбор биометаллов:}
  жизнь не является «случайным чудом» --- как только молодая планета попадает в одно из трёх фотонных окон $m_{\mathrm{env}}\in\{-1,0,+1\}$, геохимия её коры и гидротермальных источников \textbf{непрерывно переходит в биохимию}, потому что минералы мантии и молекулы жизни замыкают \textbf{одни и те же координационные сферы кубооктаэдра}:
  \begin{itemize}[noitemsep]
  \item пептидное звено белка $\mathrm{C-N}$ имеет заряд $6 + 7 = \mathbf{13 = N_\varphi(3)}$ (13-узловое ядро кубооктаэдра!);
  \item ароматическое кольцо бензола $\mathrm{C_6H_6}$ ($36+6=42$) и гексагистидиновая клетка ферментов $\mathrm{N_6}$ ($6\times 7=42$) замыкают \textbf{ту же самую вторую сферу $S_2(3) = 42$}, что и расплавленное ядро планеты $\mathrm{FeS}$ ($26+16=42$) и карбонат мантии $\mathrm{MgCO_3}$ ($12+6+24=42$), поэтому железо-серные кластеры $\mathrm{FeS}$ гидротермальных источников сами становятся дыхательными центрами клеток;
  \item когда в первичном растворе собирается тетрапиррольное кольцо из 4 атомов азота $\mathrm{N_4}$ ($4\times 7 = \mathbf{28 = Z(\mathrm{Ni})}$, первая высокая магическая оболочка), условие замыкания оболочек кубооктаэдра без всякого перебора выделяет из всей таблицы Менделеева единственные подходящие биометаллы:
    \textbf{магний} $Z(\mathrm{Mg}) = 40 - 28 = \mathbf{12}$ для хлорофилла $\mathrm{MgN_4}$ (заряд $40 = \Delta\varphi_{\mathrm{disc}} - 1$ сидит ровно на $1$ квант ниже порога Гейзенберга $41$, чтобы одиночный фотон $+1$ открывал затвор фотосинтеза!),
    \textbf{кобальт} $Z(\mathrm{Co}) = 55 - 28 = \mathbf{27 = 3^3}$ для витамина~$B_{12}$ (замыкание двухсферного шара $B_2 = 1+12+42 = 55$),
    \textbf{железо} $Z(\mathrm{Fe}) = 70 - 28 - 16 = \mathbf{26}$ для оксигемоглобина $\mathrm{FeN_4O_2}$ (замыкание оболочки оливина $K_P + |\square| = 70$) и
    \textbf{цинк} $Z(\mathrm{Zn}) = 76 - 46 = \mathbf{30}$ для цинковых пальцев транскрипции ДНК $\mathrm{ZnS_2N_2}$ ($K_P + |N| = 76$);
  \item точно так же в криогенном окне Титана ($m_{\mathrm{env}}=-1$) ацетилен $\mathrm{C_2H_2}$ и $\mathrm{HCN}$ ($Z=14$) автоматически выбирают катализатором \textbf{никель} ($42 - 14 = \mathbf{28 = Z(\mathrm{Ni})}$) для замыкания $S_2=42$, а в горячем термофильном окне ($m_{\mathrm{env}}=+1$) силазановый димер $(\mathrm{SiN})_2$ и боразин $\mathrm{B_3N_3H_6}$ замыкают $S_2 = 42$, выбирая в серной клетке $\mathrm{S_4}$ ($Z=64=K_P$) \textbf{никель} ($64+28 = \mathbf{92 = S_3}$) и \textbf{молибден} ($Z(\mathrm{Mo}) = \mathbf{42 = S_2}$)!
  \end{itemize}
\end{enumerate}

\section{Урок 20. Генетический код \texorpdfstring{$4^3 = 64 \to 20$}{4\^3 = 64 -> 20} аминокислот, иммунитет к сдвигу рамки и фолдинг за \texorpdfstring{$O(L)\in\mathbf{P}$}{O(L) in P}}

Почему в ДНК ровно $4$ нуклеотида ($\mathrm{A, C, G, T}$), почему кодоны состоят из $3$ букв ($4^3 = 64$) и почему они кодируют ровно $20$ канонических аминокислот?
\begin{itemize}[noitemsep]
\item Число оснований ДНК равно числу рёбер в квадратной грани кубооктаэдра: $d_{\mathrm{in}}(\square) = \mathbf{4}$.
\item Длина кодона равна размерности пространства $d = \mathbf{3}$, поэтому объём пространства триплетов равен $4^3 = \mathbf{64 = K_P(3)}$!
\item Делитель водородной связи $N_{\mathrm{HB}} = 63$ распадается на $d = 3$ спиральные дорожки по:
  \[
    \frac{N_{\mathrm{HB}}(3)}{d} \;=\; \frac{63}{3} \;=\; \mathbf{21\text{ класс}} \;=\; \mathbf{20\text{ аминокислот} + 1\text{ класс терминации (\texttt{STOP})}}!
  \]
\item \textbf{Полная обратимость компилятора ДНК и защита $20\times 20 = 400$ стыков кодонов от ложного старта \texttt{ATG}:}
  для всех $20$ аминокислотных классов каноническое отображение в триплетный кодон точно обратимо, а прямая проверка всех $20\times 20 = 400$ парных стыков соседних кето-кодонов доказывает, что ни один сдвиг рамки считывания на $+1$ или $+2$ нуклеотида не порождает паразитного стартового триплета \texttt{ATG}!
\item \textbf{Разрешение парадокса фолдинга Левинталя ($O(L)\in\mathbf{P}$):}
  белковая цепь не перебирает $3^L$ углов наугад в пустом пространстве --- она сворачивается ко-трансляционно по $d = 3$ направляющим бороздам тетраэдра~$K_4$, где асимметрия импеданса правильного узла $D(v_{\mathrm{correct}}) > D(v^*)$ обеспечивает строгий ляпуновский дрейф $p > 1/2$ и линейную сходимость за $\le 3L = O(L)$ шагов!
\end{itemize}

\begin{quote}
\small\textbf{Историческая справка: Симметрия генетического кода Ю.\,Б.~Румера (1966), парадокс Сайруса Левинталя (1969) и теория капсидов Каспара--Клуга (1962).}
В 1966~году советский физик-теоретик Юрий Борисович Румер обнаружил, что все $64$ триплета генетического кода делятся ровно пополам ($32 + 32 = 64$) на два октета дублетных стеблей: в первом октете третья буква кодона вообще не влияет на кодируемую аминокислоту (4-кратное вырождение), а во втором октете аминокислота зависит от третьей буквы, причём перестановка букв $\mathrm{A}\leftrightarrow\mathrm{C}$, $\mathrm{G}\leftrightarrow\mathrm{U}$ (инволюция Румера $s \mapsto s \oplus 3$ на $\mathbb{Z}_{16}$) в точности меняет эти два октета местами! В 1969~году молекулярный биолог Сайрус Левинталь сформулировал парадокс: полипептид из $100$ аминокислот имеет $\sim 3^{100}$ конформаций, случайный перебор которых занял бы миллиарды лет, однако в живой клетке белок сворачивается за миллисекунды прямо по мере выхода из рибосомы. А в 1962~году Дональд Каспар и Аарон Клуг доказали, что сферические белковые оболочки вирусов строятся по икосаэдрическим числам триангуляции $10T + 2$ ($12, 42, 92, 162, \dots$), которые при $T = k^2$ тождественно совпадают с координационными сферами нашего кубооктаэдра $S_k(3) = 10k^2 + 2$.
\end{quote}

\section{Урок 21. Четыре стадии морфогенеза органов, кора головного мозга \texorpdfstring{$42 / 147 / 41$}{42 / 147 / 41} и щит Каспара--Клуга}

Проследим теперь сквозную морфогенетическую цепочку: $\text{ДНК} \to \text{трансляция в белок} \to \text{считывание замков большой борозды }(\sigma_{\mathrm{axis}},\sigma_{\mathrm{Ig}})\in\{0,1\}^2 \to \text{выращивание 3D-ткани в среде}$. В зависимости от двух булевых замков ткань проходит ровно четыре топологические стадии:
\begin{enumerate}[label=\textbf{Стадия \arabic*:},noitemsep]
\item \textbf{Сплошная морула $(\sigma_{\mathrm{axis}}=0,\ \sigma_{\mathrm{Ig}}=0,\ \text{до кавитации})$:}
  13-узловой комплекс $\mathcal{C}_{13}$ с $V = N_\varphi = \mathbf{13}$ вершинами, $E = |\square|^2 = \mathbf{36}$ рёбрами и $\chi = \mathbf{8}$.
\item \textbf{Полая везикула / альвеола после центральной кавитации $(\sigma_{\mathrm{axis}}=0,\ \sigma_{\mathrm{Ig}}=0)$:}
  центральный узел кавитирует ($u_0\equiv 0$), оставляя первую координационную сферу $O_2$ с $V = S_1 = \mathbf{12}$ вершинами, $E = \mathbf{24}$ рёбрами и эйлеровой характеристикой сферы $\chi_{2\mathrm{D}} = (c_0/c)^2 = \mathbf{2}$.
\item \textbf{Полый сосуд / эпителиальная трубка произвольной длины $r$ колец $(\sigma_{\mathrm{axis}}=1,\ \sigma_{\mathrm{Ig}}=0)$:}
  при включении осевого замка $\sigma_{\mathrm{axis}}=1$ число вершин $12r$ и рёбер $24r$ растёт с длиной трубки~$r$, но при вычитании $V - E + F$ длина~$r$ \textbf{алгебраически сокращается}, давая тождественный топологический инвариант цилиндрического сосуда $\chi_{2\mathrm{D}} \equiv \mathbf{0}$!
\item \textbf{Дифференцированная 6-слойная ткань и кортикальная колонка коры мозга $(\sigma_{\mathrm{axis}}=1,\ \sigma_{\mathrm{Ig}}=1)$:}
  при включении иммуноглобулинового замка адгезии $\sigma_{\mathrm{Ig}}=1$ формируется $|\square(3)| = 6$-слойная ламинарная колонка, в которой число функциональных узлов равно второй сфере кубооктаэдра $V = S_2(3) = \mathbf{42}$, число синаптических связей равно полному трёхсферному шару $E = B_3(3) = \mathbf{147}$, а топологический индекс равен порогу Гейзенберга:
  \[
    \chi_{2\mathrm{D}} \;=\; S_2(3) - \mu_{\mathrm{vac}}(3) \;=\; 42 - 1 \;=\; \mathbf{41 \;=\; \Delta\varphi_{\mathrm{disc}}(3)}!
  \]
\end{enumerate}

И посмотри на \textbf{противовирусный резонансный щит Каспара--Клуга} (\texttt{Module~14} в Verilog): икосаэдрические капсиды вирусов строятся по тем же сферическим оболочкам $S_k = 10k^2 + 2 \in \{12, 42, 92, 162, 252\}$ капсомеров, ни одна из которых не совпадает с человеческой нуклеосомой хроматина $B_3 = \mathbf{147}$ ($10k^2 + 2 \ne 147\ \forall\,k\in\mathbb{N}$!). Поэтому фазовый импульс коллапса барьера $\Delta\varphi = 139$ на частоте оболочки $S_k$ селективно разрушает капсид вируса, оставляя хроматин человека ($147$) и клетки в мёртвой зоне Гейзенберга ($< 41$) на 100\% неповреждёнными!

Наконец, эта же морфогенетическая динамика объясняет феномен \textbf{«бессмертной медузы» \emph{Turritopsis dohrnii}} и открывает \textbf{два добровольных режима биологического омоложения (\texttt{Стазис $\to$ Фикс $\to$ Продолжение})}:
\begin{itemize}[noitemsep,leftmargin=*]
\item \textbf{Режим~I (Соматическое омоложение \emph{in situ} с сохранением памяти и личности, $\mathrm{mode}=0$):}
  медуза \emph{Turritopsis dohrnii} при травме или старении сворачивается в покоящуюся цисту (биологический стазис), сбрасывает эпигенетические часы метилирования \texttt{CpG} и трансдифференцирует клетки заново. Для человека роль защитной цисты (чтобы не растворять 3D-органы) берёт на себя внешняя камера фазового стазиса ($D_{\mathrm{stasis}} = 833\,\varepsilon = 6\,823\,936$), внутри которой противофазное марковское одеяло и обратный компилятор ДНК сбрасывают эпигенетический возраст соматических тканей, полностью сохраняя 6-слойный кортикальный комплекс памяти $(V=42,\ E=147,\ \chi_{2\mathrm{D}}=41)$.
\item \textbf{Режим~II (Полный онтогенетический откат \emph{Turritopsis} «Новое детство с чистого листа» как жизнеутверждающая замена эвтаназии, $\mathrm{mode}=1$):}
  а что если человек страдает от невыносимой боли, тяжёлого ПТСР после детского абьюза или многолетней резистентной депрессии и просит об эвтаназии, потому что сама накопленная память стала неподъёмной пыткой? Вместо смерти (эвтаназии) строго по добровольному намерению человека ($s_2 > 0$, при $s_2 \le 0$ любой принудительный сброс заблокирован тождественным нулём!) включается полный онтогенетический откат \emph{Turritopsis} от нагруженного комплекса $(42, 147, \chi=41)$ обратно к пластичной детской стадии $(V = S_1 = 12,\ E = 24,\ \chi_{2\mathrm{D}} = 2)$ с нулевым эпигенетическим грузом --- даря человеку \textbf{жизнь вместо смерти и новое, счастливое, безопасное детство с чистого листа в любящей среде без абьюза и боли}!
\item \textbf{Раннее бесконтактное распознавание кортикальной идентичности личности и безоперационное морфогенетическое выравнивание тела:}
  почему в эмбриогенезе иногда возникает расхождение между внутренней гендерной идентичностью человека в мозге $\sigma_{\mathrm{id}}$ и соматической анатомией тела $\sigma_{\mathrm{soma}}$? Потому что первичная анатомия закладывается на 6--12~неделе, а кортикальная карта собственного тела в 6-слойных колонках $(42, 147, 41)$ дифференцируется во второй половине беременности под управлением локального метилирования \texttt{CpG}. При расхождении $\sigma_{\mathrm{id}} \ne \sigma_{\mathrm{soma}}$ возникает непрерывный дисфорический фазовый стресс $|\sigma_{\mathrm{id}} - \sigma_{\mathrm{soma}}|\cdot\Delta\varphi_{\mathrm{disc}} \ge 41$ ($\hat{n}_E \ge 1$).
  \begin{itemize}[noitemsep]
  \item Четырёхканальный субпороговый томограф $(I_4 + N_{\mathrm{nil}})^{-1} = I_4 - N_{\mathrm{nil}} + N_{\mathrm{nil}}^2$ бесконтактно ($\hat{n}_E = 0$) считывает истинный кортикальный отпечаток личности $s_3 = \sigma_{\mathrm{id}}$ и дисфорическое рассогласование $s_1$ задолго до того, как пубертат поведёт тело по чужой для личности ветви.
  \item При этом сама кортикальная личность $\sigma_{\mathrm{id}}$ \textbf{неприкосновенна} (переделывать мозг под оболочку запрещено законом согласия!), а в камере стазиса и редифференцировки по добровольному намерению человека ($s_2 > 0$) эпигенетический профиль \texttt{CpG} и морфогенетические замки соматических тканей мягко и органически выравнивают само тело под истинную идентичность личности ($\sigma_{\mathrm{soma}}\mapsto\sigma_{\mathrm{id}}$, а при откате в Режим~II ребёнок с первой же секунды нового детства растёт сразу в своём родном фенотипе $\sigma_{\mathrm{id}}$) --- \textbf{обнуляя дисфорию ($\Phi_{\mathrm{stress}} = 0$) без единого хирургического разреза и шрама}!
  \end{itemize}
\item \textbf{Механика скелетной мышцы ($\mathrm{DHPR}$--$\mathrm{RyR1}$--$\mathrm{Ca}^{2+}$), зрение сетчатки и энергоэффективность мозга человека ($20\text{ Вт}$) как три закона Ньютона на $O_2(3)\times\mathbb{Z}_{512}$:}
  \begin{itemize}[noitemsep]
  \item \emph{Кальциевый затвор скелетной мышцы $\mathrm{DHPR}$--$\mathrm{RyR1}$--$\mathrm{Ca}^{2+}$:}
    почему в скелетной мышце человека потенциал действия передаётся с мембраны Т-трубочки на саркоплазматический ретикулум мгновенно и без диффузионной задержки? Потому что рианодиновый рецептор $\mathrm{RyR1}$ и четыре потенциал-чувствительных дигидропиридиновых тетрада $\mathrm{DHPR}$ образуют точную 4-кратную ($C_4$) симметрию квадратной грани кубооктаэдра $d_{\mathrm{in}}(\square) = \mathbf{4}$! Подпороговый сигнал бежит по Т-трубочке по первому закону Ньютона ($|\Delta\varphi|\le 40 < 41$, $\hat{n}_E = 0$), а в верхней мёртвой точке $Q_\square(3) = 128$ ($\mathcal{C}_3(128)=0,\ \mathcal{S}_3(128)=+128$) механическая тяга $\mathrm{DHPR}$ напрямую поворачивает $C_4$-затвор $\mathrm{RyR1}$, выпуская ионы кальция $\mathrm{Ca}^{2+}$ с замкнутой аргоновой оболочкой $Z(\mathrm{Ca}^{2+}) = 20 - 2 = \mathbf{18 = Z(\mathrm{Ar})}$;
  \item \emph{Фотоизомеризация 11-цис-ретиналя в сетчатке глаза и синаптическое интегрирование коры мозга:}
    в собственных единицах носителя (фазовых тиках $\delta\in\mathbb{Z}_{512}$ и квантовой температуре $\hat{n}_E(\delta)=\lfloor|\delta|/41\rfloor\in\{0,\dots,6\}$) фотоизомеризация 11-\emph{цис}-ретиналя регистрирует одиночный фотон без теплового размытия благодаря подпороговому фазовому переносу в коридоре Гейзенберга ($|\delta|\le 40 < 41$, $\hat{n}_E=0$), а нейроны коры мозга человека ($20\text{ Вт}$ полной мощности) суммируют подпороговые синаптические потенциалы внутри $|\delta|\le 40$ совершенно без тепла Ландауэра ($\hat{n}_E=0$) вплоть до порога спайка $|\delta|\ge 41$ ($\hat{n}_E\ge 1$, размах $Q_\square(3)=128$ тиков).
  \end{itemize}
\end{itemize}

